\section{一维卷积的基本操作}

\subsection{从视觉到语言：2D 卷积到 1D 卷积}

CNN 最初为计算机视觉设计，提供了平移不变性——无论目标出现在图像的哪个位置，都能被检测到。在视觉中，卷积核（filter/mask）在 2D 图像上滑动，通过点积计算每个位置的特征值。

\begin{figure}[H]
\centering
\includegraphics[width=0.95\textwidth]{slides-images/slide-007.jpg}
\caption{视觉中的 2D 卷积：卷积核在图像上滑动，计算每个位置的特征值\protect\footnotemark}
\end{figure}
\footnotetext{Slides 第8页。}

对于语言，我们只需要进行 \textbf{1D 卷积}——因为文本是一维序列。每个词有一个 $d$ 维的词向量，卷积核覆盖连续的 $h$ 个词（$h$ 即 n-gram 的大小），在序列上从头滑动到尾。

\begin{figure}[H]
\centering
\includegraphics[width=0.95\textwidth]{slides-images/slide-008.jpg}
\caption{1D 卷积用于文本：卷积核（红色部分）在词向量序列上滑动，计算每个 trigram 位置的特征值\protect\footnotemark}
\end{figure}
\footnotetext{Slides 第9页。}

\subsection{单滤波器卷积的计算过程}

对于一个 trigram 滤波器，设词向量维度为 $d=4$，滤波器参数为 $\mathbf{w} \in \mathbb{R}^{3 \times 4}$。对于位置 $i$ 处的 trigram（词 $i$, $i+1$, $i+2$），计算过程为：

$$c_i = f(\mathbf{w} \cdot \mathbf{x}_{i:i+2} + b)$$

其中：
\begin{itemize}
    \item $\mathbf{x}_{i:i+2}$ 是三个词向量的拼接
    \item $\mathbf{w}$ 是滤波器权重（与位置无关，参数共享）
    \item $b$ 是偏置项
    \item $f$ 是非线性激活函数（如 sigmoid、ReLU）
\end{itemize}

\begin{knowledgebox}{卷积的参数共享特性}
与全连接层不同，卷积层的核心特性是\textbf{参数共享}：同一个滤波器在所有位置使用相同的权重。这不仅大幅减少了参数量，还赋予了模型平移不变性——无论特定的 n-gram 模式出现在句子的哪个位置，都能被同样的滤波器检测到。
\end{knowledgebox}

\subsection{Padding（填充）}

在不使用填充的情况下，对长度为 $n$ 的句子使用大小为 $h$ 的滤波器，输出长度为 $n - h + 1$（缩短了）。为了保持输出长度与输入一致，可以使用 \textbf{zero padding}：

\begin{itemize}
    \item \textbf{Same padding}：在两端各填充 $\lfloor h/2 \rfloor$ 个零向量，输出长度 = 输入长度
    \item \textbf{Wide convolution}：填充更多零向量，输出长度 $>$ 输入长度
    \item \textbf{No padding（narrow convolution）}：不填充，输出长度 $<$ 输入长度
\end{itemize}

\begin{figure}[H]
\centering
\includegraphics[width=0.95\textwidth]{slides-images/slide-009.jpg}
\caption{Padding 示意：通过在句子两端填充零向量，使卷积输出的长度保持不变\protect\footnotemark}
\end{figure}
\footnotetext{Slides 第10页。}

\subsection{多滤波器卷积}

一个滤波器只能检测一种模式。实际应用中，我们会使用 \textbf{多个滤波器}（channels），每个滤波器学习检测不同的特征：

\begin{figure}[H]
\centering
\includegraphics[width=0.95\textwidth]{slides-images/slide-010.jpg}
\caption{多通道 1D 卷积：3 个滤波器分别对每个 trigram 位置计算特征值，得到 3 维的特征向量\protect\footnotemark}
\end{figure}
\footnotetext{Slides 第11页。}

设有 $m$ 个大小为 $h$ 的滤波器，输入序列长度为 $n$（带 padding），则：
\begin{itemize}
    \item 每个滤波器产生长度为 $n$ 的特征图（feature map）
    \item $m$ 个滤波器共同产生 $m \times n$ 的特征矩阵
    \item 每个位置有一个 $m$ 维的新表示向量
\end{itemize}

\subsection{本章小结}

1D 卷积通过在词向量序列上滑动固定大小的滤波器来提取局部特征。通过 padding 控制输出长度，通过多滤波器捕获不同类型的模式。卷积操作的参数共享特性使其高效且具有平移不变性。

%% ========== 池化与分类 ==========

